\documentclass[../../main.tex]{subfiles}% 注意这里的文件路径不能用 ./main.tex ,否则用latexmk编译子文件会报错
\graphicspath{{\subfix{./image/}}} % 指定图片目录,后续可以直接使用图片文件名
% 注意这里的文件路径不能用 ../../image/ ,否则用latexmk编译子文件会报错

% 例如:
% \begin{figure}[H]
% \centering
% \includegraphics[scale=0.3]{图.png}
% \caption{}
% \label{figure:图}
% \end{figure}
% 注意:上述\label{}一定要放在\caption{}之后,否则引用图片序号会只会显示??.

\begin{document}

\section{曲面上正交参数曲线网的存在性}

\begin{lemma}
设$f(u,v),g(u,v)$是定义在区域$D\subset\mathbb{E}^2$上的两个不同时为零的连续可微函数，则对任意一点$(u_0,v_0)\in D$，存在它的邻域$U\subset D$，以及$U$上的非零连续函数$\lambda(u,v)$，使得$\lambda(u,v)$是一次微分式
\begin{align}
\omega = f(u,v)\mathrm{d}u + g(u,v)\mathrm{d}v
\label{eq:onediffform}
\end{align}
的积分因子，即存在$U$上的连续可微函数$k(u,v)$，满足
\begin{align*}
\lambda(u,v)\big(f(u,v)\mathrm{d}u + g(u,v)\mathrm{d}v\big) = \mathrm{d}k(u,v).
\end{align*}
\end{lemma}
\begin{remark}
该引理仅对二元一次微分式成立，因此本节下述结论仅适用于曲面情形。
\end{remark}

\begin{theorem}
设正则参数曲面$S:\bm r=\bm r(u,v)$上存在两个处处线性无关的连续可微切向量场$\bm a(u,v),\bm b(u,v)$，则对任意$p\in S$，存在$p$的邻域$U\subset S$及$U$上的新参数系$(\tilde u,\tilde v)$，使得新参数曲线的切向量$\bm r_{\tilde u},\bm r_{\tilde v}$分别与$\bm a,\bm b$平行，即$\bm r_{\tilde u}\parallel\bm a,\;\bm r_{\tilde v}\parallel\bm b$。
\end{theorem}
\begin{remark}
该定理表明曲面上存在局部参数系，使参数曲线切向与给定线性无关切向量场平行，但一般无法使$\bm r_{\tilde u}=\bm a,\bm r_{\tilde v}=\bm b$。
\end{remark}
\begin{proof}

设切向量场在自然基底$\{\bm r_u,\bm r_v\}$下的表达式为
\begin{align}
\bm a(u,v) &= a_1(u,v)\bm r_u + a_2(u,v)\bm r_v,\nonumber\\
\bm b(u,v) &= b_1(u,v)\bm r_u + b_2(u,v)\bm r_v,
\label{eq:tangentfielddecomp}
\end{align}
其中$a_i,b_i\;(i=1,2)$为连续可微函数，且
\begin{align}
A=\begin{vmatrix}a_1(u,v)&a_2(u,v)\\b_1(u,v)&b_2(u,v)\end{vmatrix}\neq0.
\label{eq:detnonzero}
\end{align}

若存在容许参数变换$u=u(\tilde u,\tilde v),\;v=v(\tilde u,\tilde v)$，使得$\bm r_{\tilde u}\parallel\bm a,\;\bm r_{\tilde v}\parallel\bm b$，则存在函数$\lambda,\mu$满足
\begin{align}
\bm r_{\tilde u}=\lambda\bm a,\quad \bm r_{\tilde v}=\mu\bm b.
\label{eq:paralleltangentvec}
\end{align}
代入\eqref{eq:tangentfielddecomp}得
\begin{align*}
\bm r_{\tilde u}&=\lambda a_1\bm r_u+\lambda a_2\bm r_v,\\
\bm r_{\tilde v}&=\mu b_1\bm r_u+\mu b_2\bm r_v.
\end{align*}
结合参数变换的Jacobi矩阵定义，有
\begin{align}
J=\begin{pmatrix}
\frac{\partial u}{\partial\tilde u}&\frac{\partial v}{\partial\tilde u}\\
\frac{\partial u}{\partial\tilde v}&\frac{\partial v}{\partial\tilde v}
\end{pmatrix}
=\begin{pmatrix}
\lambda a_1&\lambda a_2\\
\mu b_1&\mu b_2
\end{pmatrix}.
\label{eq:jacobiorig}
\end{align}

由反函数的Jacobi矩阵为原矩阵的逆，可得
\begin{align}
\begin{pmatrix}
\frac{\partial\tilde u}{\partial u}&\frac{\partial\tilde v}{\partial u}\\
\frac{\partial\tilde u}{\partial v}&\frac{\partial\tilde v}{\partial v}
\end{pmatrix}
=\frac{1}{\lambda\mu A}\begin{pmatrix}
\mu b_2&-\lambda a_2\\
-\mu b_1&\lambda a_1
\end{pmatrix}.
\label{eq:jacobiinv}
\end{align}
因此
\begin{align}
\mathrm{d}\tilde u&=\frac{1}{\lambda A}\big(b_2\mathrm{d}u - b_1\mathrm{d}v\big),\nonumber\\
\mathrm{d}\tilde v&=\frac{1}{\mu A}\big(-a_2\mathrm{d}u + a_1\mathrm{d}v\big).
\label{eq:diffformtildeuv}
\end{align}
令一次微分式
\begin{align*}
\alpha = b_2\mathrm{d}u - b_1\mathrm{d}v,\quad \beta = -a_2\mathrm{d}u + a_1\mathrm{d}v.
\end{align*}
由引理，存在非零连续可微函数$\xi,\eta$分别为$\alpha,\beta$的积分因子，即
\begin{align}
\mathrm{d}\tilde u &= \xi\big(b_2\mathrm{d}u - b_1\mathrm{d}v\big),\nonumber\\
\mathrm{d}\tilde v &= \eta\big(-a_2\mathrm{d}u + a_1\mathrm{d}v\big).
\label{eq:intfactordiff}
\end{align}
此时
\begin{align*}
\begin{pmatrix}
\frac{\partial\tilde u}{\partial u}&\frac{\partial\tilde v}{\partial u}\\
\frac{\partial\tilde u}{\partial v}&\frac{\partial\tilde v}{\partial v}
\end{pmatrix}
=\begin{pmatrix}
\xi b_2&-\eta a_2\\
-\xi b_1&\eta a_1
\end{pmatrix},
\end{align*}
其行列式$\xi\eta(a_1b_2-a_2b_1)\neq0$，故$(\tilde u,\tilde v)$为$U$上的新参数系。

进一步计算切向量：
\begin{align*}
\bm r_{\tilde u}&=\bm r_u\frac{\partial u}{\partial\tilde u}+\bm r_v\frac{\partial v}{\partial\tilde u}
=\frac{1}{\xi(a_1b_2-a_2b_1)}\big(a_1\bm r_u+a_2\bm r_v\big)
=\frac{1}{\xi A}\bm a,\\
\bm r_{\tilde v}&=\bm r_u\frac{\partial u}{\partial\tilde v}+\bm r_v\frac{\partial v}{\partial\tilde v}
=\frac{1}{\eta(a_1b_2-a_2b_1)}\big(b_1\bm r_u+b_2\bm r_v\big)
=\frac{1}{\eta A}\bm b.
\end{align*}
因此$\bm r_{\tilde u}\parallel\bm a,\;\bm r_{\tilde v}\parallel\bm b$.

\end{proof}

\begin{theorem}
对正则参数曲面$S:\bm r=\bm r(u,v)$上任意一点$p\in S$，存在$p$的邻域$U\subset S$及$U$上的新参数系$(\tilde u,\tilde v)$，使得新参数曲线的切向量$\bm r_{\tilde u},\bm r_{\tilde v}$彼此正交，即$(\tilde u,\tilde v)$为$S$在$U$上的\textbf{正交参数系}。
\end{theorem}

\begin{remark}
该定理为存在性定理，构造正交参数系等价于求解对应一次微分式的积分因子；该定理保证了正则曲面局部正交参数曲线网的存在性，简化了曲面理论的局部分析。
\end{remark}

\begin{proof}

记曲面第一基本量$E=\bm r_u\cdot\bm r_u,\;F=\bm r_u\cdot\bm r_v,\;G=\bm r_v\cdot\bm r_v$，对自然基底$\{\bm r_u,\bm r_v\}$作Schmidt正交化，令
\begin{align}
\bm e_1=\frac{\bm r_u}{|\bm r_u|}=\frac{\bm r_u}{\sqrt{E}}.
\label{eq:orthobase1}
\end{align}
取$\bm b=\bm r_v+\lambda\bm e_1$，满足$\bm b\cdot\bm e_1=0$，代入得
\begin{align*}
\lambda = -\bm r_v\cdot\bm e_1 = -\frac{F}{\sqrt{E}},
\end{align*}
因此
\begin{align*}
\bm b=\bm r_v-\frac{F}{E}\bm r_u,
\end{align*}
计算模长
\begin{align*}
|\bm b|^2=G-\frac{F^2}{E}=\frac{EG-F^2}{E}.
\end{align*}
令
\begin{align}
\bm e_2=\frac{\bm b}{|\bm b|}=\frac{1}{\sqrt{EG-F^2}}\left(-\frac{F}{\sqrt{E}}\bm r_u+\sqrt{E}\bm r_v\right),
\label{eq:orthobase2}
\end{align}
则$\{\bm e_1,\bm e_2\}$为$S$上的单位正交切标架场。

由前述定理，存在局部参数系$(\tilde u,\tilde v)$使得$\bm r_{\tilde u}\parallel\bm e_1,\bm r_{\tilde v}\parallel\bm e_2$，故$\bm r_{\tilde u}\perp\bm r_{\tilde v}$，即$(\tilde u,\tilde v)$为正交参数系.

\end{proof}

\begin{example}
设空间曲线$\bm{r}=\bm{r}(s)$以弧长$s$为参数，曲率是$\kappa$。写出它的切线面的参数方程，使得相应的参数曲线构成正交曲线网。
\end{example}

\begin{example}
改写曲面
\begin{align*}
\bm{r}=(u\cos v,u\sin v,u+v).
\end{align*}
的参数方程，使得它的参数曲线网是正交曲线网。
\end{example}

\begin{example}
求曲面
\begin{align*}
\bm{r}=(v\cos u - k\sin u,v\sin u + k\cos u,ku).
\end{align*}
的参数曲线的正交轨线，其中$k>0$是常数。
\end{example}



\end{document}